\documentclass[10pt]{article}

\usepackage{microtype,amsmath,amssymb,amsthm,booktabs,array}
\usepackage[margin=0.88in]{geometry}
\usepackage[colorlinks=true,linkcolor=blue,urlcolor=blue,citecolor=blue]{hyperref}
\usepackage{bookmark}
\newtheorem{theorem}{Theorem}
\newtheorem{proposition}[theorem]{Proposition}
\newtheorem{lemma}[theorem]{Lemma}
\newtheorem{corollary}[theorem]{Corollary}
\newcommand{\dd}{\mathrm d}
\newcommand{\Tr}{\operatorname{Tr}}
\newcommand{\sech}{\operatorname{sech}}
\newcommand{\artanh}{\operatorname{artanh}}
\newcommand{\ph}{\varphi}
\newcommand{\RR}{\mathbb R}
\setlength{\emergencystretch}{2em}
\title{Uniform Bulk Modular Asymptotics\\and the Leading Wilson Entropy Law}
\author{Kirtus Justus (Theoretical Physicist)\\\small Port Coquitlam, British Columbia, Canada}
\date{October 2026}
\hypersetup{pdftitle={Uniform Bulk Modular Asymptotics and the Leading Wilson Entropy Law},pdfauthor={Kirtus Justus (Theoretical Physicist)}}
\usepackage{fontspec,unicode-math,mathtools}
\setmainfont{DejaVu Serif}\setsansfont{DejaVu Sans}\setmonofont{DejaVu Sans Mono}\setmathfont{Latin Modern Math}
\begin{document}
\maketitle
\begin{abstract}
The previously derived modular period coefficients can be connected to the
finite concentration spectrum with a uniform remainder on every fixed
interior interval of the scaled Jacobi spectrum. An exact positive
realization of the omitted differential spectrum becomes, after a second
change of variable, a smooth whole-line equation with potential bounded
away from zero. Its recessive logarithmic derivative gives an exact
integral for the concentration logit. An integrable Riccati residual
then proves an asymptotic expansion to every fixed finite order.
Deformation of the coefficient contours identifies its terms with the
original even WKB periods. In particular the expansion through $N^{-5}$
has a uniform $O(N^{-7})$ remainder for $|\lambda/N|\le\tau<1/2$.
A conservative explicit constant is supplied for $\tau=1/4$.
A separate global linear envelope then bounds the positive-majorant-weighted
modular discrepancy uniformly. Combined with the established Wilson
covariance and boundary-entropy estimates, this gives the fixed-scaled-mass
relative entropy $\mu^2\log L/6+O_\mu(1)$. Neither a limiting finite part
nor a uniform spectral edge expansion is obtained.
\end{abstract}

\section{Statement and conventions}

Let $N\ge2$ be even. For $0\le j,k<N$, define
\begin{equation}
B_{jk}=\begin{cases}1/2,&j=k,\\
\sin[\pi(j-k)/2]/[\pi(j-k)],&j\ne k,\end{cases}
\qquad J_{j-1,j}=J_{j,j-1}=\frac{j(N-j)}N,\quad J_{jj}=0.
\label{eq:model}
\end{equation}
These matrices commute, $J$ has simple spectrum, and $0<B<I$.
For their common modes write
\begin{equation}
Jv=\lambda v,\quad Bv=\zeta v,\qquad
t=\lambda/N,\quad \varepsilon=\log\frac{1-\zeta}{\zeta},\quad
e=\varepsilon+N\eta(t),
\label{eq:conventions}
\end{equation}
where
\begin{equation}
\eta(t)=\pi t\,{}_3F_2\!\left(\tfrac14,\tfrac12,\tfrac34;1,\tfrac32;4t^2\right),
\qquad |t|<1/2.
\label{eq:eta}
\end{equation}
The alternating-sign involution sends $(t,\zeta,\varepsilon,e)$ to
$(-t,1-\zeta,-\varepsilon,-e)$. All logarithms are natural.

\begin{theorem}[Uniform interior expansion]\label{thm:main}
For every $0<\tau<1/2$ and integer $K\ge1$, there are finite constants
$C_{\tau,K}$ and $N_{\tau,K}$ such that every common mode of
\eqref{eq:model} with $|t|\le\tau$ and even $N\ge N_{\tau,K}$ satisfies
\begin{equation}
\left|e-\sum_{j=1}^K\frac{\ph_{2j-1}(t)}{N^{2j-1}}\right|
\le C_{\tau,K}N^{-2K-1}.
\label{eq:main}
\end{equation}
Here the $\ph_{2j-1}$ are the even WKB period coefficients defined in
Section~\ref{sec:periods}. They are real analytic and odd on $(-1/2,1/2)$.
For $K=3$ they are precisely the previously reduced functions
\begin{equation}
\ph_1=\tfrac t2\eta'+(t^2-\tfrac16)\eta'',\qquad
\ph_3=A_3\eta'+B_3\eta'',\qquad
\ph_5=A_5\eta'+B_5\eta'',
\label{eq:closed}
\end{equation}
with the rational coefficients in \eqref{eq:AB3}--\eqref{eq:B5}.
For $|t|\le1/4$ one admissible, deliberately conservative bound is
\begin{equation}
\left|e-\frac{\ph_1(t)}N-\frac{\ph_3(t)}{N^3}-\frac{\ph_5(t)}{N^5}\right|
\le \frac{2273823860\,\pi}{N^7},\qquad N\ge16.
\label{eq:explicit}
\end{equation}
\end{theorem}

This proves the fixed-interior spectral statement left open in
\cite{periods}. The bounds apply also across $t=0$; the apparent
singularities there in \eqref{eq:closed} are removable. The constants
depend on the distance from $|t|=1/2$. No convergence of the full
inverse-power series is asserted. The exact identities used from
\cite{tails,moments} are stated next, so the dependence on those
earlier results is explicit.

\section{Exact spectral input and a whole-line resolvent}

On $I=(-\pi/2,\pi/2)$ let $p(x)=\cos x$ and let
\[
A_N=-(p u')'-\frac{N^2-1}{4}p u
\]
have its Friedrichs realization. Its simple eigenvalues are
$\mu_{N,0}<\mu_{N,1}<\cdots$. The first $N$ eigenvalues are
$-\operatorname{spec}(NJ/2)$. The remaining spectrum is positive,
and the exact tail formula of \cite{tails} is
\begin{align}
F_N(b)&=\sum_{k\ge N}\log\frac{\mu_{N,k}+b}{\mu_{N,k}-b}
       =2\sum_{k\ge N}\artanh\frac b{\mu_{N,k}},\label{eq:tail}\\
F_N(N\lambda/2)&=\log\frac\zeta{1-\zeta}.\label{eq:mode}
\end{align}
The series converges locally uniformly for $|b|<\mu_{N,N}$.
In \cite{moments} the same tail is realized exactly by
\begin{equation}
T_N=-D_z^2+\frac{4N^2-1}{16p(z)}+\frac{p(z)}{16},\qquad
\frac{\dd z}{\dd x}=p^{-1/2},
\label{eq:positive}
\end{equation}
with its Friedrichs realization on the finite $z$ interval.
In particular $\mu_{N,N}>N^2/4$ and $T_N^{-1}$ is trace class.

These are operator and boundary-condition statements, not just formal
differential substitutions. Briefly, the finite Fourier subspace
exhausts the first $N$ modes by a nodal-count argument.
The logarithmic derivative of the endpoint connection coefficient is
the resolvent trace of $A_N$. Removing its finite spectral polynomial
and integrating that trace from $-b$ to $b$ gives \eqref{eq:tail}--\eqref{eq:mode},
without an unspecified entire factor. Successive ground-state
Darboux deletions realize the omitted spectrum. The transformed
Fourier Wronskian is a nonzero constant times $p^{N^2/4}$, yielding
\eqref{eq:positive}. At the initial critical endpoint the
Friedrichs condition excludes the logarithmic branch; each deletion
raises the selected exponent by one. The final endpoints are limit
point. These domain arguments are proved in \cite{tails,moments};
the algebraic deletion itself belongs to the classical framework
of Crum~\cite{crum}.

Undoing the Liouville transformation in \eqref{eq:positive} gives
the unitarily equivalent operator on $L^2(I,\dd x)$
\begin{equation}
\widetilde T_N=-(p u')'+\left(\frac{N^2}{4p}+\frac p4\right)u.
\label{eq:tilde}
\end{equation}
Indeed its Liouville correction is
$p_{xx}/4-p_x^2/(16p)=-3p/16-1/(16p)$.

Introduce the second coordinate
\begin{equation}
r=\log(\sec x+\tan x),\qquad p=\sech r,\qquad \dd x=p\,\dd r,
\qquad h=N^{-1}.
\label{eq:coordinate}
\end{equation}
For $b=N^2t/2$, the quadratic form of $\widetilde T_N-b$ becomes
\begin{equation}
\int_{\RR}\left\{|u_r|^2+
 \left(\frac{N^2}4+\frac{p^2}4-bp\right)|u|^2\right\}\dd r.
\label{eq:form}
\end{equation}
It acts in $L^2(\RR,p\,\dd r)$ and its form domain is $H^1(\RR)$
when $|t|\le\tau<1/2$. Define on ordinary $L^2(\RR,\dd r)$
\begin{equation}
P_h(t)=-h^2D_r^2+Q_t+h^2v,\qquad
Q_t=\frac{1-2t\sech r}{4},\qquad v=\frac{\sech^2r}{4}.
\label{eq:wholeline}
\end{equation}
Put $s_0=\sqrt{1-2\tau}/2$. Then $Q_t\ge s_0^2>0$.
All derivatives of $Q_t-1/4$ and $v$ decay exponentially, uniformly
for $|t|\le\tau$.

\begin{lemma}[Weighted resolvent identity]\label{lem:resolvent}
Let $Uu=\sqrt p\,u$ map $L^2(p\,\dd r)$ unitarily to $L^2(\dd r)$.
Then
\begin{equation}
U(\widetilde T_N-N^2t/2)^{-1}U^{-1}
=h^2p^{1/2}P_h(t)^{-1}p^{1/2}.
\label{eq:weightedresolvent}
\end{equation}
If $G_h(t;r,s)$ is the Green kernel of $P_h(t)$ and
$\mathcal F_N(t)=F_N(N^2t/2)$, then
\begin{equation}
\mathcal F_N'(t)=\frac12\int_{\RR}p(r)
 [G_h(t;r,r)+G_h(-t;r,r)]\,\dd r.
\label{eq:Fprime}
\end{equation}
\end{lemma}
\begin{proof}
The identity follows first on compactly supported tests from
\eqref{eq:form}: the weak equation is $N^2P_h(t)u=p f$.
Coercivity and density extend it to the closed forms and their
bounded inverses. This is a weighted resolvent identity; it does
not identify the ordinary whole-line spectrum of $P_h(t)$ with
the discrete tail spectrum.
Positivity and the diagonal bound proved below give
$0<G_h(t;r,r)\le(2hs_0)^{-1}$. Since $\int_{\RR}p\,\dd r=\pi$,
the weighted diagonal is integrable. On compact intervals the
positive continuous kernel has trace equal to its diagonal integral;
exhaustion gives the same identity on the line. Equivalently,
trace class also follows directly from the left side of
\eqref{eq:weightedresolvent}. Differentiating \eqref{eq:tail} gives
$\mathcal F_N'=(N^2/2)\Tr[(T_N-b)^{-1}+(T_N+b)^{-1}]$.
Substitute \eqref{eq:weightedresolvent} and $N^2h^2=1$.
\end{proof}

\section{An exact Riccati integral for the logit}

Let $u_-(r;t,h)$ be the solution of $P_h(t)u=0$ normalized by
$u_-(r)\sim e^{r/(2h)}$ as $r\to-\infty$. Set
\begin{equation}
w=h\frac{u_-'}{u_-},\qquad
hw'+w^2=Q_t+h^2v,\qquad w(-\infty)=1/2.
\label{eq:riccati}
\end{equation}
The potential is strictly positive. For precision, $w$ can be
constructed by solving \eqref{eq:riccati} from $r=-R$ with value
$1/2$ and letting $R\to\infty$. Lower and upper constant barriers
give $s_0\le w\le\max(1/2,\sqrt{\sup(Q_t+h^2v)})$.
The difference of two such solutions satisfies a first-order
equation with damping at least $2s_0/h$, giving convergence on
compact sets and uniqueness of the bounded solution with the
declared left limit. Comparing with $1/2$ gives
\begin{equation}
h(w-1/2)'+(w+1/2)(w-1/2)=-tp/2+h^2v.
\label{eq:integrable}
\end{equation}
The right side is integrable. The integrating-factor formula
therefore shows $w-1/2\in L^1(\RR)$, uniformly on compact $t$
intervals for fixed $h$, and $w\to1/2$ at both ends.
Reconstruction by integration fixes the required normalization of
$u_-$ and proves that it is positive. Define
\begin{equation}
\mathcal A_h(t)=\int_{\RR}(w(r;t,h)-1/2)\,\dd r.
\label{eq:A}
\end{equation}

\begin{proposition}[Exact connection formula]\label{prop:exact}
For $|t|<1/2$,
\begin{equation}
\boxed{\quad\mathcal F_N(t)=
\frac{\mathcal A_h(-t)-\mathcal A_h(t)}h.\quad}
\label{eq:exactRiccati}
\end{equation}
\end{proposition}
\begin{proof}
Difference quotients in $t$ satisfy a damped first-order equation
with source $-p/2$. Its integrating-factor formula is dominated
in $L^1$ by an exponential convolution of $p$.
Taking the limit justifies differentiation and gives
\begin{equation}
\partial_t w(r)=-\frac1{2h}\int_{-\infty}^r p(s)
 \exp\!\left[-\frac2h\int_s^r w(a)\,\dd a\right]\dd s.
\label{eq:derivativew}
\end{equation}
Reduction of order constructs the right recessive solution as
$u_-(r)\int_r^\infty u_-(a)^{-2}\dd a$. Its Wronskian with
$u_-$ is $-1$, so
\begin{equation}
G_h(t;s,s)=\frac1{h^2}\int_s^\infty
 \exp\!\left[-\frac2h\int_s^r w(a)\,\dd a\right]\dd r
\le\frac1{2hs_0}.
\label{eq:green}
\end{equation}
Absolute integrability, or Tonelli applied to the nonnegative
integrands, permits integration of \eqref{eq:derivativew} in $r$:
\[
\mathcal A_h'(t)=-\frac h2\int_{\RR}p(s)G_h(t;s,s)\,\dd s.
\]
The derivative of the right side of \eqref{eq:exactRiccati} is
therefore \eqref{eq:Fprime}. Both sides vanish at $t=0$.
This determines the integration constant exactly.
\end{proof}

\section{A uniform residual estimate}

Smooth positive-potential Liouville--Green expansions and their
error bounds are classical \cite{dlmf,dunster}. Here an elementary
$L^1$ estimate suffices and keeps the constants tied directly to
the exact logit.

For fixed $K$, recursively choose $a_j(r,t)$ so that
\begin{align}
S_K&=\sqrt{Q_t}\left(1+\sum_{j=1}^Kh^{2j}a_j\right),\label{eq:S}\\
\mathcal R_K&=S_K^2-Q_t-h^2v
 -\frac{h^2}2(\log S_K)''
 +\frac{h^2}4[(\log S_K)']^2\label{eq:residual}
\end{align}
has zero coefficients through $h^{2K}$. Primes in this section
mean $r$ derivatives. At order $h^{2j}$ the new unknown occurs
only as $2Q_t a_j$, so the recursion is unique. Every $a_j$ is
rational in $p=\sech r$ and $t$, with possible denominators only
powers of $1-2tp$. Moreover $a_j(0,t)=0$ when regarded as a
function of $p$. These facts follow inductively using
\begin{equation}
D_r^2f=p^2(1-p^2)f_{pp}+p(1-2p^2)f_p,
\qquad (D_rf)(D_rg)=p^2(1-p^2)f_pg_p.
\label{eq:derivativealgebra}
\end{equation}
Thus $a_j$ and every fixed $r$ derivative are $O(p)$ uniformly
for $|t|\le\tau$. In particular all integrals below converge.

Choose $h_0>0$ so small that, for $0<h\le h_0$,
\begin{equation}
\sum_{j=1}^Kh^{2j}\|a_j\|_\infty\le\tfrac12,
\qquad
h\left(\tfrac12\|Q_t'/Q_t\|_\infty+
       2\sum_{j=1}^Kh^{2j}\|a_j'\|_\infty\right)\le s_0/2.
\label{eq:hsmall}
\end{equation}
Here the norms include the supremum over $|t|\le\tau$.
They are finite, so such an $h_0$ exists. Then $S_K\ge s_0/2$ and
\begin{equation}
w_K=S_K-\frac h2(\log S_K)'\ge s_0/4,
\qquad hw_K'+w_K^2-(Q_t+h^2v)=\mathcal R_K.
\label{eq:wK}
\end{equation}

\begin{lemma}[Integrated stability bound]\label{lem:stability}
Under \eqref{eq:hsmall},
\begin{equation}
\|w-w_K\|_1\le s_0^{-1}\|\mathcal R_K\|_1.
\label{eq:stability}
\end{equation}
Consequently
\begin{align}
\left|\mathcal F_N(t)-\frac1h\int_{\RR}
 [S_K(r;-t,h)-S_K(r;t,h)]\,\dd r\right|
\le\frac{\|\mathcal R_K(\cdot,t,h)\|_1+
             \|\mathcal R_K(\cdot,-t,h)\|_1}{hs_0}.
\label{eq:certificate}
\end{align}
\end{lemma}
\begin{proof}
For $d=w-w_K$ the exact equation is
\[
hd'+(w+w_K)d=-\mathcal R_K,\qquad d(-\infty)=0.
\]
The integrating factor and $w+w_K\ge s_0$ yield
\[
|d(r)|\le\frac1h\int_{-\infty}^r
 |\mathcal R_K(s)|e^{-s_0(r-s)/h}\,\dd s.
\]
Integrating gives \eqref{eq:stability}. Also $S_K\to1/2$ at both
ends, so $\int_{\RR}(\log S_K)'\dd r=0$. Substitute
\eqref{eq:wK} into \eqref{eq:A} and use \eqref{eq:exactRiccati}.
\end{proof}

The cancellation defining $a_j$ is sufficient for an actual
uniform integrable residual, not merely a pointwise formal one.
Indeed put $z=h^2$ and $A=1+\sum_{j=1}^K z^ja_j$.
If
\[
R_0=-\tfrac14(\log Q_t)''+\tfrac1{16}[(\log Q_t)']^2,
\]
then direct calculation gives the finite polynomial identity
\begin{equation}
A^2\mathcal R_K=Q_t(A^4-A^2)+z(R_0-v)A^2
 -\frac z2 AA''+\frac{3z}4(A')^2
 +\frac z4(\log Q_t)'AA'.
\label{eq:cleared}
\end{equation}
Its coefficients through $z^K$ vanish. Each remaining coefficient
is $O(p)$, uniformly on $0\le p\le1$, $|t|\le\tau$; this follows
from \eqref{eq:derivativealgebra} and $a_j(0,t)=0$.
The denominators are bounded away from zero. Since $A\ge1/2$,
there is a finite $M_{\tau,K}$, given by these finitely many
rational functions, for which
\begin{equation}
|\mathcal R_K(r,t,h)|\le M_{\tau,K}p(r)h^{2K+2},
\quad \|\mathcal R_K\|_1\le\pi M_{\tau,K}h^{2K+2}.
\label{eq:L1residual}
\end{equation}
Equation~\eqref{eq:certificate} now gives a uniform $O(h^{2K+1})$
logit error. The residual formula also provides a direct finite-$h$
bound whenever its norms and the positivity conditions are bounded.

Define
\begin{equation}
\Phi_{2j-1}(t)=\int_{\RR}
 [\sqrt{Q_t}\,a_j(r,t)-\sqrt{Q_{-t}}\,a_j(r,-t)]\,\dd r.
\label{eq:realcoeff}
\end{equation}
The uniformly convergent binomial expansion and
$\int_{\RR}\sech^m r\,\dd r=
\sqrt\pi\,\Gamma(m/2)/\Gamma((m+1)/2)$ give
\begin{equation}
\int_{\RR}[\sqrt{Q_{-t}}-\sqrt{Q_t}]\,\dd r=\eta(t).
\label{eq:leading}
\end{equation}
It follows that the exact spectral interpolation
$\mathfrak e_N(t)=N\eta(t)-\mathcal F_N(t)$ obeys
\begin{equation}
\mathfrak e_N(t)=\sum_{j=1}^K\Phi_{2j-1}(t)N^{1-2j}
 +O_{\tau,K}(N^{-2K-1})
\label{eq:preperiod}
\end{equation}
on the whole interval $[-\tau,\tau]$. The integrals in
\eqref{eq:realcoeff} are odd and real analytic there. At a common
finite mode, \eqref{eq:mode} identifies $\mathfrak e_N(t)$ with $e$.

\section{Identification with the original periods}\label{sec:periods}

For completeness, the finite Fourier equation for an eigenmode is
\[
(\cos\xi\,U')'+\left[\frac{N^2-1}{4}\cos\xi-\frac{N\lambda}2\right]U=0.
\]
Writing $u=\cos\xi$ and $V=\sqrt u\,U$ gives
\begin{equation}
V''+[N^2q(u,t)+r_0(u)]V=0,\qquad
q=\frac{u-2t}{4u},\qquad r_0=\frac1{4u^2}.
\label{eq:originalode}
\end{equation}
The formal even momentum is
\begin{equation}
W=N\sqrt q\left(1+\sum_{j\ge1}N^{-2j}f_j(u,t)\right),\qquad
W^2+\tfrac12(\log W)_{\xi\xi}
 -\tfrac14[(\log W)_\xi]^2=N^2q+r_0.
\label{eq:originalmomentum}
\end{equation}
For $0<|t|<1/2$, let $\Gamma$ be counterclockwise around the cut
joining $0$ and $2t$, excluding $\pm1$, and put
\begin{equation}
\ph_{2j-1}=i\oint_\Gamma\sqrt q\,f_j\,\dd\xi,
\qquad \dd\xi=-\frac{\dd u}{\sqrt{1-u^2}}.
\label{eq:perioddef}
\end{equation}
The branch of $\sqrt{1-u^2}$ is positive on $(-1,1)$ with cuts
on the outer real rays; $\sqrt q\sim1/2$ at infinity and has its
cut between $0$ and $2t$. These conventions agree with \cite{periods}.

\begin{proposition}[Exact equality of coefficient integrals]\label{prop:periods}
For every $j\ge1$ and $|t|<1/2$,
$\Phi_{2j-1}(t)=\ph_{2j-1}(t)$, with the latter extended at zero.
\end{proposition}
\begin{proof}
In \eqref{eq:originalmomentum} set $\xi=ir$, so that
$u=\cosh r=1/p$. Second derivatives and products of first
derivatives change sign, while $q=Q_t$ and $r_0=v$.
The resulting even-momentum equation is exactly
\eqref{eq:residual} with its residual set formally to zero.
Uniqueness of the coefficient recursion gives
\begin{equation}
a_j(p,t)=f_j(1/p,t).
\label{eq:inversion}
\end{equation}
In particular $f_j(u,t)=O(1/u)$ at infinity. It is rational,
with possible poles only at $0$ and $2t$, and it satisfies
$f_j(-u,t)=f_j(u,-t)$.

For $j\ge1$ the differential in \eqref{eq:perioddef} is
$O(u^{-2})\dd u$ at infinity. Thus $\Gamma$ can be deformed
to the two outer branch cuts: the circle at infinity contributes
zero, and the small arcs about $\pm1$ vanish because their
singularities have order $1/2$. On the upper bank of the positive
outer cut, $\sqrt{1-u^2}=-i\sqrt{u^2-1}$; on the upper bank of the
negative cut it is $+i\sqrt{u^2-1}$. The lower values have opposite
signs. With the stated counterclockwise orientation this yields
\begin{equation}
\ph_{2j-1}=2\int_1^\infty
 \frac{\sqrt{q(u,t)}f_j(u,t)-\sqrt{q(u,-t)}f_j(u,-t)}
 {\sqrt{u^2-1}}\,\dd u.
\label{eq:outercuts}
\end{equation}
Substitute $u=\cosh r$ and use evenness in $r$ to obtain
\eqref{eq:realcoeff}. The analytic continuation through $t=0$
follows from that uniformly convergent real integral.
The leading coefficient is handled separately by
\eqref{eq:leading}; the unsubtracted leading differential
does not have the decay needed for this contour deformation.
\end{proof}

Together, \eqref{eq:preperiod}, Proposition~\ref{prop:periods}
and \eqref{eq:mode} prove \eqref{eq:main}.
For the first three functions the exact rational reductions of
\cite{periods} give \eqref{eq:closed}, with
\begin{align}
A_3&=\frac{108t^4-83t^2+7}{120t(4t^2-1)^2},&
B_3&=\frac{708t^6-789t^4+153t^2-14}{360t^2(4t^2-1)^2},\label{eq:AB3}\\
A_5&=-\frac{\begin{gathered}6384t^{10}-51320t^8+26567t^6\\[-2pt]
-13594t^4+2019t^2-124\end{gathered}}{1680t^3(4t^2-1)^4},\label{eq:A5}\\
B_5&=-\frac{\begin{gathered}38304t^{12}-374688t^{10}+267522t^8\\[-2pt]
-170573t^6+39472t^4-4937t^2+248\end{gathered}}{5040t^4(4t^2-1)^4}.
\label{eq:B5}
\end{align}
These formulae were proved by exact derivatives on the period
curve; the present argument supplies their spectral connection
and remainder. Their small-$t$ behavior includes
$\ph_1(t)=\pi t/4+O(t^3)$, consistent with the separately proved
plunge-law sign $e\simeq-\varepsilon/(4N^2)$ in this convention.
The absolute uniform bound above alone is not a relative error
bound for arbitrarily small nonzero $t$.

\begin{corollary}[Interior spectral projection]
Let $\Pi_{N,\tau}=\mathbf1_{[-\tau,\tau]}(J/N)$ and
$E_N=\log(I-B)-\log B+N\eta(J/N)$. Under the conditions of
Theorem~\ref{thm:main},
\[
\left\|\Pi_{N,\tau}\left[E_N-
 \sum_{j=1}^KN^{1-2j}\ph_{2j-1}(J/N)\right]\Pi_{N,\tau}\right\|
\le C_{\tau,K}N^{-2K-1}.
\]
\end{corollary}
\begin{proof}
Apply \eqref{eq:main} in the common orthonormal eigenbasis.
\end{proof}

\input{Weighted_Consequences.tex}

\section{Numerical checks and remaining scope}

The proof does not use fitted constants or numerical mode values.
As a separate check of its exact normalization, set $z=e^r$ and
$u_-=z^{N/2}f(z)$. Equation~\eqref{eq:wholeline} becomes
\begin{equation}
z(1+z^2)^2f''+(N+1)(1+z^2)^2f'
 +[N^2t(1+z^2)-z]f=0,\qquad f(0)=1.
\label{eq:numericalODE}
\end{equation}
Frobenius expansion followed by overlapping Taylor expansions
computes $f(1),f'(1)$. Evenness of the potential gives the growing
connection coefficient
\[
\mathfrak a_h(t)=\frac{2f(1)[Nf(1)/2+f'(1)]}{N},\qquad
\mathcal F_N(t)=\log\mathfrak a_h(-t)-\log\mathfrak a_h(t).
\]
Here $\log\mathfrak a_h=\mathcal A_h/h$, as follows by integrating
$w-1/2$ over the line. This computation uses the whole-line
equation, whereas the retained reference values were computed
from finite Jacobi eigenvectors and concentration quadratic forms.
For four $N=160$ modes, at $t$ approximately
$-0.04909,-0.12511,-0.25111,-0.40225$, the absolute differences
in the logits are respectively
$3.90\times10^{-89}$, $7.61\times10^{-90}$,
$2.27\times10^{-88}$ and $4.78\times10^{-88}$.
This is consistent with the retained input precision. Separately, direct
110-digit finite-matrix calculations at every mode for
$N=2,4,8,16,32,64$ satisfy both global logit bounds, the discrepancy
envelope, and the particle--hole sign convention.

Exact symbolic substitution verifies \eqref{eq:inversion} for
$j=1,2,3$ and the four vanishing coefficients of
\eqref{eq:cleared} below $h^8$. Real quadrature of
\eqref{eq:realcoeff} at six points between $-0.3$ and $0.45$
agrees with \eqref{eq:closed} to absolute errors below
$6\times10^{-66}$. A further check evaluates the exact
interpolation at fixed $t$; these points need not be finite
Jacobi eigenvalues. Table~\ref{tab:remainder} displays the scaled
remainder after the first three terms.

\begin{table}[ht]
\centering
\begin{tabular}{rrrr}
\toprule
$N$ & $t=0.125$ & $t=0.3$ & $t=0.45$\\
\midrule
16  & $-2.756381$ & $-286.9294$ & $-2.020035\times10^6$\\
32  & $-2.854713$ & $-327.4185$ & $-4.221313\times10^6$\\
64  & $-2.881067$ & $-340.4583$ & $-6.240675\times10^6$\\
128 & $-2.887778$ & $-343.9819$ & $-7.261610\times10^6$\\
\bottomrule
\end{tabular}
\caption{$N^7[\mathfrak e_N(t)-\ph_1(t)/N-\ph_3(t)/N^3-\ph_5(t)/N^5]$
from the exact whole-line interpolation. The much larger values
near $t=1/2$ illustrate the importance of fixing the interior interval.}
\label{tab:remainder}
\end{table}

All decimal checks use multiprecision arithmetic with increased
precision and truncation for replication; they are not outward-rounded
interval enclosures. The constant in \eqref{eq:explicit}, by contrast,
comes from exact rational inequalities in the appendix. It is much
larger than the observed errors and is not presented as a sharp bound.

The result resolves the fixed-interior period connection and
remainder problem for the half-filled sine kernel. It leaves open
uniform control in a moving edge window $1/2-|t|\to0$.
Consequently the earlier coefficient-level Stirling limits do not
yet become a uniform spectral edge theorem. The bound on the
interior projection does not by itself control $\|E_N\|$ on the full
spectrum. The separate global envelope and Wilson spectral tail
suffice for the leading entropy law, but do not establish a limiting
finite part or the proposed inverse-logarithmic rate.
General filling and the continuum prolate problem require separate
exact operators and domain analysis; they are not covered by
Theorem~\ref{thm:main}. Within \emph{From Prephysical Hypotheses to
Physical Models}, these are dimensionless reference-model results.
They neither select a physical clock nor establish a universal
physical model or a quantum-field-theoretic mass gap.

\appendix
\section{A finite rational certificate for the explicit constant}

Here $K=3$, $|t|\le1/4$, $p\in[0,1]$, $z=h^2$ and
$A=1+za+z^2b+z^3c$. The first three coefficients can be generated by
\begin{align*}
R_0&=-\tfrac14(\log Q)''+\tfrac1{16}[(\log Q)']^2,
&a&=(v-R_0)/(2Q),\\
R_1&=-\tfrac12a''+\tfrac14(\log Q)'a',
&b&=-(Qa^2+R_1)/(2Q),\\
L_2&=b-a^2/2,\qquad
R_2=-\tfrac12L_2''+\tfrac14(\log Q)'L_2'+\tfrac14(a')^2,
&c&=-(2Qab+R_2)/(2Q).
\end{align*}
Each $a_j=n_j(p,t)/(1-2pt)^{3j}$ for a polynomial $n_j$
divisible by $p$. For example
\[
a=\frac{p(p^3t^2-pt^2+p-2t)}{2(1-2pt)^3}.
\]
For a polynomial $n=\sum n_{ab}p^at^b$, the rational number
$\sum|n_{ab}|4^{-b}$ bounds $|n|$ on this whole rectangle.
Applying this to $n_j/p$, $n_j$, and its first two $p$ derivatives,
and using $1-2pt\ge1/2$, gives the following uniform majorants:
\begin{center}
\begin{tabular}{crrr}
\toprule
$j$ & $|a_j|/p\le M_j$ & $|D_ra_j|/p\le D_j$
 & $|D_r^2a_j|/p\le E_j$\\
\midrule
1 & $13/2$ & $31$ & $379/2$\\
2 & $7353/8$ & $17595/2$ & $386379/4$\\
3 & $6923583/16$ & $49698847/8$ & $1566168253/16$\\
\bottomrule
\end{tabular}
\end{center}
For example, with $m=3j$, $\delta=1/2$ and the indicated
polynomial norms,
\begin{align*}
\|\partial_pa_j\|_\infty
&\le\|n_j'\|\delta^{-m}+2m\tau\|n_j\|\delta^{-m-1},\\
\|\partial_p^2a_j\|_\infty
&\le\|n_j''\|\delta^{-m}+4m\tau\|n_j'\|\delta^{-m-1}
 +4m(m+1)\tau^2\|n_j\|\delta^{-m-2}.
\end{align*}
Equation~\eqref{eq:derivativealgebra} then bounds $D_j$ by the
first expression and $E_j$ by the sum of these two expressions.

Take $h_0=1/16$ and replace $s_0$ conservatively by $s_*=1/3$.
The two left sides of \eqref{eq:hsmall} are bounded, respectively, by
\[
\frac{17504063}{268435456}<\frac12,\qquad
\frac{117523487}{1073741824}<\frac16=\frac{s_*}2.
\]
Thus $A\ge1/2$ and the stability argument applies for every
$h\le1/16$.

To specify the residual bound without a long expanded polynomial,
let $M(z)=1+\sum_{j=1}^3M_jz^j$,
$D(z)=\sum_{j=1}^3D_jz^j$, and $E(z)=\sum_{j=1}^3E_jz^j$.
Write $[z^k]$ for coefficient extraction. Since
$Q\le3/8$, $|R_0-v|/p\le13/16$, and
$|(\log Q)'|/p\le1$, the coefficient of $z^k$ in the
right side of \eqref{eq:cleared}, for $4\le k\le12$, has
absolute value at most $p C_k$, where
\begin{equation}
C_k=[z^k]\left\{\tfrac38(M^4+M^2)
 +z\left[\tfrac{13}{16}M^2+\tfrac12ME
                    +\tfrac34D^2+\tfrac14MD\right]\right\}.
\label{eq:Ck}
\end{equation}
Every product in these positive-power coefficients contains a
factor vanishing at $p=0$, so the single factor $p$ is retained.
The coefficients for $k<4$ vanish exactly by the recurrence.
Since $A^{-2}\le4$,
\[
\|\mathcal R_3\|_1\le4\pi h^8
 \sum_{k=4}^{12}C_k h_0^{2k-8}.
\]
The two signs in \eqref{eq:certificate} and $s_*=1/3$ give
a multiplier $24\pi\sum_{k=4}^{12}C_k16^{-2k+8}$.
Exact rational evaluation of this finite expression yields
\[
24\sum_{k=4}^{12}C_k16^{-2k+8}<2273823860.
\]
This proves \eqref{eq:explicit} using finite rational inequalities.

\paragraph{Acknowledgments and computational assistance.}
ChatGPT (OpenAI) assisted with the present analysis, symbolic
identities, numerical implementations and proof scrutiny. The
earlier endpoint connection proposal was supplied through work
with Claude (Anthropic); the precise identities and assumptions
used here are recorded in \cite{tails,moments}. The broader program
has also used Gemini (Google), Grok (SpaceXAI) and DeepSeek
(High-Flyer). This does not attribute the present checks to those
systems or imply endorsement. The mathematical argument and
reproduction materials remain open to independent review;
AI-assisted checking is not independent human peer review.

\begin{thebibliography}{9}
\bibitem{periods} K. Justus, \emph{Higher-Order Modular Periods for the
Half-Filled Sine Kernel}, research manuscript, October 2026.
\bibitem{tails} K. Justus, \emph{Spectral Tails and Exact Concentration
Logits}, research manuscript, October 2026.
\bibitem{moments} K. Justus, \emph{Positive Moment Bounds and Growing
Plunge Windows}, research manuscript, October 2026.
\bibitem{crum} M. M. Crum, Associated Sturm--Liouville systems,
\emph{Quart. J. Math. Oxford (2)} \textbf{6} (1955), 121--127.
\bibitem{dlmf} NIST Digital Library of Mathematical Functions,
\S2.7(iii), Liouville--Green (WKBJ) Approximation.
\url{https://dlmf.nist.gov/2.7#iii}.
\bibitem{dunster} T. M. Dunster, \emph{Liouville--Green expansions
of exponential form, with an application to modified Bessel functions},
2017. \href{https://arxiv.org/abs/1708.00548}{arXiv:1708.00548}.
\bibitem{majorants} K. Justus, \emph{Exact Spectral Majorants for a
Wilson--Dirac Interval}, research manuscript, 23 October 2026.
\bibitem{entropy} K. Justus, \emph{Relative Entropy of Finite-Interval
Wilson--Dirac Quasifree States}, research manuscript, 23 October 2026;
the bounded nonlinear, boundary-entropy and leading-order reduction
theorems, distinct from the conditional finite-part theorem.
\bibitem{positivetrace} K. Justus, \emph{Positive Reference-Modular
Trace for Wilson--Dirac Intervals}, research manuscript, October 2026;
the independent finite trace reduction, not its later asymptotic consequence.
\end{thebibliography}
\end{document}
